%!TEX root = TQFTnotes.tex
%%%%%%%%%%%%%%%%%%%%%%%%%%%%
\chapter*{Preface}\label{c:preface}
This text is based on the notes for a course taught by the
author at Stony Brook University in the Spring of 2026.

This is a work in progress, to be updated frequently.


The goal of this course is to give a review of mathematics related
to Topological Quantum Field Theory and more generally, extended TQFT ---
from its origin in physics through the initial definition of Atiyah to notions
of extended TQFT following works of Lurie and possibly to recent
developments related to factorization homology. We also
include constructions of TQFT via state sums (also known as lattice
TQFT), at least in examples.

Our goal is to make it a readable overview; therefore, in many places we only
give a sketch of the proof, or even of a definition (for example, we do not give
all details of construction of $(\infty, n)$ categories),  referring the
reader to original papers for details. Instead, we concentrate on examples of
TQFTs, such as
\begin{itemize}
  \item Gauge theory with finite group $G$ (also known as Dijkgraaf--Witten
  theory)
  \item 2d TQFT using symmetric Frobenius algebras, which we present both using
  generators for the category of 2d bordisms (Schommer-Pries)
  and using state sums
  (Fukuma--Hosono--Kawai theory)
  \item Turaev--Viro 3d TQFT
  \item Crane--Yetter 4d TQFT and the closely related Chern--Simons 3d TQFT
\end{itemize}

Of all the examples mentioned above, Chern--Simons theory is by far the best
known, and there are numerous expositions of this theory. Partly for this
reason, we do not plan to spend too much time on it; in particular, we will not
be discussing quantum invariants of knots and links.




\clearpage
